\section{案例簇 B：Autonomous Driving 与交叉口协调}

Stone 复盘了自动驾驶研究中的一个代表性问题：当道路系统主体都变成 autonomous agent，交通规则会不会从静态信号灯转向协商式 reservation system。

\begin{knowledgebox}{Reservation-based intersection 的 Agent 含义}
每辆车都可视为 agent，进入路口前向调度系统请求时空轨迹槽位（reservation）：
\begin{itemize}
    \item 获批后按保序轨迹通过；
    \item 未获批则等待；
    \item 目标是减少全局等待和冲突概率。
\end{itemize}
\end{knowledgebox}

这个想法对今天 LLM Agent 也有直接启发：在共享资源系统中，``协作协议'' 往往比单 Agent 局部最优更重要。

\begin{warningbox}{混合交通是难点}
纯 autonomous agent 环境容易优化；human + autonomous 混合环境会引入策略不确定性和行为非平稳性，这是部署中更难的阶段。
\end{warningbox}

\subsection{本章小结}
自动驾驶案例强调了一个事实：Agent intelligence 不是离散任务成功率，而是协议化协同能力。

